
\section{Présentation du système}

Vous travaillez sur une cellule d'\textbf{ouverture et de vidage de cartons}. Elle est utilisée pour vider automatiquement des cartons contenant un produit connu (par exemple, des sachets) sur un convoyeur spécialisé.

Le système est composé de plusieurs sous-systèmes :
\begin{itemize}
    \item \textbf{Un convoyeur d'alimentation} qui achemine les cartons fermés vers la zone de traitement.
    \item \textbf{Un robot 6 axes} muni d'une pince de préhension. 
    \item \textbf{Ladite pince de préhension} : 
    \begin{itemize}
        \item Permet de saisir un carton,
        \item Est équipée de deux paires de bloqueurs permettant de maintenir les volets du carton ouverts.
    \end{itemize}
    \item \textbf{Un poste de découpe}, équipé d’un actionneur de lame, qui découpe le ruban adhésif fermant le carton.
    \item \textbf{Un poste d'ouverture} permettant de soulever les volets du carton à l’aide de vérins.
    \item \textbf{Un poste de vidage} où le carton est retourné et vidé.
    \item \textbf{Un convoyeur de vidage} qui transporte le contenu du carton vers une station aval.
\end{itemize}

Le système fonctionne de manière cyclique selon les étapes suivantes :
\begin{enumerate}
    \item Le robot saisit un carton sur le convoyeur d'alimentation.
    \item Il le présente au poste de découpe pour sectionner le ruban adhésif.
    \item Il le présente ensuite au poste d'ouverture :
    \begin{itemize}
        \item Les grands volets sont soulevés par le poste d'ouverture,
        \item Le robot les maintient ouverts avec ses bloqueurs,
        \item Le robot effectue une rotation pour exposer les petits volets,
        \item Ceux-ci sont également soulevés, puis maintenus ouverts.
    \end{itemize}
    \item Le carton est ensuite amené au poste de vidage :
    \begin{itemize}
        \item Une plaque métallique se positionne devant l’ouverture du carton,
        \item Le robot retourne le carton, le contenu étant retenu par la plaque,
        \item Il effectue une translation pour déverser le contenu sur le convoyeur,
        \item Un vérin pneumatique pousse le produit restant sur le convoyeur.
    \end{itemize}
\end{enumerate}

\section*{Objectif du devoir}

Ce devoir porte sur la \textbf{modélisation orientée objet} et la \textbf{programmation} du système de préhension. Il se concentre particulièrement sur :

Vous utiliserez les conventions de nommage et les structures proposées à l’IUT. Le raisonnement objet et la capacité à structurer proprement votre code seront des éléments d’évaluation essentiels.
